\begin{table}[H]
    \centering
    \fontsize{10}{12.5}\selectfont
    \caption{Resultados computacionales de distancias y tiempos óptimos calculados desde Plaza Mayor}
    \label{tab:resultados_distancias}
    \begin{tabularx}{\textwidth}{clccX}
        \toprule
        \textbf{Nodo} & \textbf{Nombre de la Intersección} & \textbf{Distancia Pura ($m$)} & \textbf{Tiempo Efectivo ($s$)} & \textbf{Secuencia de Ruta Óptima} \\
        \midrule
        \textbf{A} & Plaza Mayor del Cusco & 0.0 & 0.0 & \texttt{A} \\
        \textbf{B} & Plaza Regocijo (Cusipata) & 180.0 & 44.8 & \texttt{A -> B} \\
        \textbf{C} & Qorikancha (Santo Domingo) & 450.0 & 108.0 & \texttt{A -> C} \\
        \textbf{D} & Plaza San Francisco & 400.0 & 101.6 & \texttt{A -> B -> D} \\
        \textbf{E} & Mercado Central San Pedro & 650.0 & 165.3 & \texttt{A -> B -> D -> E} \\
        \textbf{F} & Plazoleta Limacpampa & 800.0 & 144.4 & \texttt{A -> C -> F} \\
        \bottomrule
    \end{tabularx}
    \vspace{0.15cm}
    \raggedright
    \textit{Nota.} Solución calculada desde el vértice fuente $s = \text{A}$ (Plaza Mayor).
\end{table}
